\documentclass[exercice,a4]{thedocument}%
\usepackage[french]{babel}%
\usepackage{pstricks}%
\startdatakeys
    \source{21GENMATPO1}
    \BO{2026}
    \cycle{4}
    \competencesDuSocle{CA,RA}
    \objectifsApprentissage{C4B3ca}
    \motsCles{Pythagore,aire,trigonométrie,angle,périmètre,Thalès}
\enddatakeys
\begin{document}%

\medskip

\parbox{0.6\linewidth}{Sur la figure ci-contre, qui n'est pas en vraie grandeur, le point C est le point d'intersection des droites (BE) et (AD).

\medskip

\begin{enumerate}
\item Démontrer que le triangle ABC est rectangle en C.
\item Calculer l'aire du triangle ABC.
\item Calculer une valeur approchée au degré près de l'angle $\widehat{\text{BAC}}$.
\item Calculer le périmètre du triangle CDE.
\item Les droites (AB) et (DE) sont-elles parallèles ?
\end{enumerate}
}\hfill
\parbox{0.35\linewidth}{
\psset{unit=0.75cm}
\begin{pspicture}(-3.3,-4.5)(2.5,5.8)
\pspolygon[linewidth=1.25pt](-3,0)(0,5.3)(0,-4.2)(2.4,0)%ABED
\uput[l](-3,0){A} \uput[u](0,5.3){B} \uput[d](0,-4.2){E} \uput[r](2.4,0){D} \uput[dl](0,0){C}
\uput[u](-1.5,0){8 cm} \rput{90}(-0.3,-2.1){12~cm}\rput{60}(1.2,-2.6){13~cm}
\rput{90}(0.3,2.65){15~cm}\rput{60}(-1.5,3.15){17~cm}
\end{pspicture}
}

\bigskip

\showthesource
\end{document}
